\begin{textsample}{NEG Dim 2 – se – Score 1.00 – se/se\_ef\_18.txt}  \label{ex:f2_neg_264}
Definição de problemas - Observar o mundo a sua volta e fazer perguntas. - Analisar demandas, delinear problemas e planejar investigações.
Propor hipóteses.
Levantamento, análise e representação - Planejar e realizar atividades de campo ( experimentos, observações, leituras, visitas, ambientes virtuais etc. ).
Desenvolver e utilizar ferramentas, inclusive digitais, para coleta, análise e representação de dados ( imagens, esquemas, tabelas, gráficos, quadros, diagramas, mapas, modelos, representações de sistemas, fluxogramas, mapas conceituais, simulações, aplicativos etc. ). - Avaliar informação ( validade, coerência e adequação ao problema formulado ). - Elaborar explicações e/ou modelos. - Associar explicações e/ou modelos à evolução histórica dos conhecimentos científicos envolvidos.
Selecionar e construir argumentos com base em evidências, modelos e/ou conhecimentos científicos.
Aprimorar seus saberes e incorporar, gradualmente, e de modo significativo, o conhecimento científico - Desenvolver soluções para problemas cotidianos usando diferentes ferramentas, inclusive digitais. - Organizar e/ou extrapolar conclusões.
Comunicação - Relatar informações de forma oral, escrita ou multimodal. - Apresentar, de forma sistemática, dados e resultados de investigações. - Participar de discussões de caráter científico com colegas, professores, familiares e comunidade em geral.
Considerar contra-argumentos para rever processos investigativos e conclusões.
Intervenção - Implementar soluções e avaliar sua eficácia para resolver problemas cotidianos. - Desenvolver ações de intervenção para melhorar a qualidade de vida individual, coletiva e socioambiental.
Fonte : BNCC, 2017

% matched lemmas: 
\end{textsample}
